\documentclass[11pt]{article}
\usepackage[margin=0.78in]{geometry}
\usepackage{amsmath,amssymb,bm}
\usepackage{booktabs}
\usepackage{enumitem}
\usepackage{hyperref}
\hypersetup{colorlinks=true, linkcolor=black, urlcolor=black}
\setlength{\parindent}{0pt}
\setlength{\parskip}{0.55em}

\title{Cyrus Guided Notes: MIT 6.245 Multivariable Control Systems}
\author{MIT course pack}
\date{2026-05-15}

\begin{document}
\maketitle

\section*{Course source}
\textbf{Official source:} \url{https://ocw.mit.edu/courses/6-245-multivariable-control-systems-spring-2004/}

\textbf{Why this course is in the system.} Graduate control course for MIMO systems, robustness, and frequency-domain design.

\textbf{Learning output.} Keep it in the graduate control chapter instead of forcing it into the first beginner round.

\section*{Prerequisite checkpoint}
\begin{itemize}[leftmargin=1.4em]
  \item State space
  \item Frequency response
  \item Linear systems
\end{itemize}

\section*{Formula checkpoint}
Do not memorize these first. Read each formula as a sentence: what changes, what is observed, what is optimized, and what output it produces.
\begin{align*}
  \bm{y}(s)=G(s)\bm{u}(s)\\
  \bar{\sigma}(G(j\omega))\\
  \|T_{zw}\|_{\infty}<\gamma\\
  S=(I+GK)^{-1},\quad T=GK(I+GK)^{-1}
\end{align*}

\section*{Visual page}
Draw MIMO as many inputs pushing many outputs, then show singular values as worst-direction gain.

\section*{GoodNotes pages}
\begin{itemize}[leftmargin=1.4em]
  \item Page MIT-M001: MIMO transfer matrix
  \item Page MIT-M002: sensitivity and robustness
\end{itemize}

\section*{Web interaction}
A gain-direction widget that rotates an input vector and shows output amplification.

\section*{Obsidian and Notion}
\begin{tabular}{p{0.22\linewidth}p{0.68\linewidth}}
\toprule
Layer & Output \\
\midrule
Obsidian & MIT Control -> 6.245 multivariable control \\
Notion & Robust-control expansion queue. \\
\bottomrule
\end{tabular}

\section*{First study move}
Open the official source, copy the prerequisite checkpoint into GoodNotes, then write one formula in your own words before watching or reading more.

\end{document}
